\cotefile{156-8}
\scanpage{90}\mspage{81}
\uline{1 juillet}\qquad \unsure{Magma associé|Magmas associés} à un ensemble ordonné $L$.

\deleted{Je vais \unsure{d}}

A) Supposons d'abord $L$ \uline{totalement} \uline{ordonné}, et \uline{divisible} i.e.
\[
 \forall a,b\in L,\quad a<b\Longrightarrow
 \underbrace{]a,b[}_{\substack{\shortparallel\,\text{déf}\\
     \{x\unsure{\in|\ni}L\mid a<x<b\}}}\ne\varnothing
\]

Dans ce cas, on va définir un \deleted{magma} \uline{\unsure{ens multistr.}}, ayant $L$ comme ens de points (ou lieux), et un
\[
 M\subset\deleted{\unsure{\mathcal P(L)|\mathrm{P}(L)}}\,\unsure{\text{Figél}|\text{Figel}}(L)
\]
défini ainsi\,:
\[
\left\{\begin{aligned}
 &M=M_0\amalg M_1,\quad\text{avec}\\
 &M_0=\big\{\deleted{\unsure{\{\{x\}\}}}\inserted{P_x}\mid x\in L\big\},\quad
     \text{où}\quad \deleted{\unsure{X}}P_x=\{\{x\}\}\\
 &M_1=\{\deleted{\unsure{X}}S_{a,b}\mid a,b\in L,\ a<b\},\quad
     \deleted{\unsure{X}}S_{ab}=\{[a,b],\{a\},\{b\}\}
\end{aligned}\right.
\]
% slanted
\margin{\uline{NB} Quand on écrit $S_{a,b}$, on sous-entend $a<b$. On a besoin de la divisibilité, pour que $S_{a,b}$ soit bien une figure, i.e. \deleted{$S_{a,b}^{\circ}=\unsure{]a,b[}$} $S_{ab}^{\circ}\overset{\text{déf}}{=}|S_{a,b}|\setminus|\partial S_{ab}|\ne\varnothing$}

Il est tautologique que les conditions $\text{Magm }1,2$ sont vérifiées. ($M$ contient les $\{\{x\}\}$ i.e. les $\unsure{X_x|P_x}$, et $M$ stable par passage à une sous-\deleted{\unsure{multi-str}}multistrate.)

Explicitons les relations $\leq,\ \overset{\circ}{\ll},\ |\bullet|$

% slanted
\margin{NB Il vaut mieux utiliser la notation $X\triangleleft Y$ pour $X<Y$ et $X\trianglelefteq Y$ pour $X\leq Y$ (notations d'une relation d'incidence) pour réserver la notation $X<Y$ à la \unsure{rel} ``positionnelle'' $X\underset{\mathrm{pos}}{\leq}Y$. cf p. 97,98.}

\[
\begin{aligned}
1^\circ)\quad X\leq Y\Longleftrightarrow X=Y\ \text{ou}\quad
  X=P_x,Y=S_{a,b},\quad&\text{avec}\quad x\in\{a,b\}\\[-2pt]
 &\text{i.e.}\quad x=a\ \text{ou}\ x=b
\end{aligned}
\]
\[
\deleted{2^\circ)\quad X\overset{\circ}{\ll}Y\quad\text{sss}\quad X=Y\ \text{ou}\ X=P_x,Y}
\]

En d'autres termes, on a
\[
\left\{\begin{aligned}
 \partial P_x&=\varnothing\\
 \partial S_{a,b}&=\{P_a,P_b\}
\end{aligned}\right.
\]

\[
\begin{aligned}
2^\circ)\quad X\overset{\circ}{\ll}Y\Longleftrightarrow{}&
  \deleted{X=Y\ \text{ou}\ \unsure{(}X=P_a,}\\[-2pt]
&\left\{\begin{aligned}
 a)&\ \text{si }X=P_x,Y=P_y:\quad x=y\quad\text{i.e. }P_x=P_y\\
 b)&\ \text{si }X=P_x,\unsure{Y|X}=S_{ab}:\quad a<x<b\quad
       [\text{i.e. }x\in|X|^\circ]\\
 c)&\ \text{si }\unsure{-}X=S_{a,b},\ Y=S_{a',b'}:\quad a'\leq a<b\leq b'\\
   &\qquad[\text{i.e.}\ \unsure{-}\ |X|\subset|Y|,\ \inserted{\text{\deleted{ou} aussi}}\,
                 |X|^\circ\subset|Y|^\circ]
 \end{aligned}\right.
\end{aligned}
\]

On voit donc\({}^{*}\) que les $P_x$ sont les éléments minimaux pour $\leq$, \uline{et} les éléments minimaux pour $\overset{\circ}{\ll}$ (grâce à divisibilité), ce sont donc aussi les él. minimaux pour $\overset{*}{\ll}$ ou lieux
% slanted
\margin{\({}^{*}\) c'est d'ailleurs tout vu, en vertu des \unsure{seules|seuls} a\unsure{.}\deleted{\unsure{xiô}}mes $\text{Magm }1,2$.}
% archive sheet number: 90

